\section{IV. THẢO LUẬN}

\subsection{4.1. Bình luận kết quả nghiên cứu \#1}
Kết quả phân tích tiền khả thi khẳng định tiềm năng vượt trội của việc ứng dụng Process Mining so với các chỉ số thống kê mô tả tĩnh truyền thống. Sự tiến bộ học tập không chỉ đơn thuần được đo bằng tổng thời lượng trực tuyến hay số lượng tài nguyên đã xem, mà được phản ánh chân thực thông qua cách thức sinh viên di chuyển qua các hoạt động học tập theo thời gian.

Việc tính toán chỉ số phù hợp quy trình theo từng cửa sổ thời gian tạo ra một góc nhìn đa chiều về nỗ lực và chiến lược học tập của sinh viên. Hướng đi này giúp giảng viên nhận diện chính xác sự tiến bộ của người học thay vì chỉ nhìn vào điểm số cuối kỳ.

\subsection{4.2. Bình luận kết quả nghiên cứu \#2}
Về phương diện lý luận, Động cơ gợi ý lộ trình dựa trên Process Mining thể hiện ưu thế rõ rệt so với các phương pháp gợi ý dựa trên nội dung (Content-based filtering). Bằng cách khai thác vị trí chính xác của người học trong chuỗi quy trình, hệ thống có thể phát hiện đúng điểm nghẽn (chẳng hạn việc sinh viên làm đi làm lại một bài tập mà không xem video hướng dẫn) để đưa ra phản hồi chính xác.

Bên cạnh đó, việc rà soát kỹ thuật cho thấy các công cụ mã nguồn mở hiện nay hoàn toàn đáp ứng được yêu cầu phân tích dữ liệu quy trình, giúp giảm thiểu rào cản chi phí khi triển khai.

\subsection{4.3. Bình luận kết quả nghiên cứu \#3}
Do nghiên cứu hiện dừng lại ở bước đánh giá tiền khả thi trên dữ liệu thứ cấp, một số giả thuyết khoa học cần tiếp tục được kiểm chứng ở giai đoạn thực nghiệm đầy đủ:
\begin{enumerate}
    \item Mức độ tương quan thực tế giữa Điểm Tiến bộ Chuỗi (LPS) và các thước đo kết quả học tập độc lập.
    \item Giá trị thông tin bổ sung của chuỗi thứ tự hành vi so với các chỉ số tần suất tương tác truyền thống.
    \item Tác động hai mặt của cơ chế cảnh báo tự động đối với động lực học tập của nhóm sinh viên có kết quả thấp.
\end{enumerate}

\subsection{4.4. Hạn chế nghiên cứu và định hướng nghiên cứu}
Nghiên cứu tiền khả thi tồn tại hai hạn chế chính. Thứ nhất, nhật ký sự kiện trên Hệ thống Quản lý Học tập chỉ ghi nhận được các tương tác trực tuyến, chưa phản ánh được các hoạt động tự học nhóm hoặc học tập ngoại tuyến. Thứ hai, chi phí tính toán của thuật toán đối sánh chuỗi có xu hướng tăng cao khi mở rộng quy mô phân tích.

Về định hướng tiếp theo, nếu thỏa mãn các ngưỡng quyết định Go/No-go, nghiên cứu sẽ triển khai thực nghiệm quy mô nhỏ trên một học phần kéo dài một học kỳ ($n \approx 200$ sinh viên), phối hợp với giảng viên xây dựng Mô hình quy trình tham chiếu và kiểm định toàn diện mô hình đánh giá.
